\begin{abstract}
    \noindent
高质量的 LLM 请求调度需要同时实现两个关键目标：一是被路由到的实例是否拥有可加速请求执行的 {\kvcache}，二是不同实例之间的工作负载是否保持均衡。
同时满足这两个目标并不容易，因为追求其中一个目标时，往往会削弱另一个目标。
现有方法通常采用各种组合器（例如线性组合），将分别反映这两个目标的指标合成为一个调度分数；但这类方法较为复杂，要么需要大量面向具体工作负载的超参数调优，要么需要开发同时感知模型与硬件的模拟器，而且即便如此仍可能得到次优性能。

本文表明，只需将两个经过精心选择的指标相乘作为调度分数——一个用于感知 {\kvcache}（若请求被路由到某实例时需要新生成的预填充 token 数），另一个用于感知负载均衡（实例当前的批大小）——就能在无需任何超参数调优的情况下，同时较好地实现这两个目标。
其核心思想是：乘法分数以一种类似线性组合的方式同时考虑两个目标，并且在实例间比较时具有一个很好的性质，即原本的超参数会在比较过程中被抵消，因此我们无需调参来寻找最佳参数。
这两个指标的选择来自我们对 LLM 特性的分析；大量实验表明，与 vLLM-v1 和某生产环境调度器相比，这种简单方法在覆盖聊天机器人、API 调用和代码智能体等真实工作负载上，可将 TTFT 分别降低 92\% 和 52\%，并将 TPOT 分别降低 21\% 和 20\%。
我们还从数学上推导了乘法可能失效的条件，并发现这类条件在实践中极为罕见，而且可以事先被检测出来（并加以缓解）。

\end{abstract}
